\section{Near-field aerodynamic corrections}
\label{sec:nearfield}

\Cref{mod:disk} treats an isolated rotor in an unbounded fluid at rest. A
machine hovering at head height in a room violates all three conditions: there
is a floor below, a ceiling above, other rotors alongside, and after a short
time the air is not at rest but circulating. This section quantifies the first
three and measures the fourth.

\subsection{Ground effect}

When the wake is constrained by a surface a distance $z$ below the rotor, the
streamtube cannot contract freely, the induced velocity at the disk falls, and
the rotor produces more thrust for the same power.

\begin{model}[Cheeseman-Bennett]
\label{mod:ige}
Modelling the ground by an image source, Cheeseman and
Bennett~\cite{cheeseman1955} obtain, at constant power,
\begin{equation}
  \frac{T_{\mathrm{IGE}}}{T_{\mathrm{OGE}}}\bigg|_{P}
  = \frac{1}{1-\left(\dfrac{R}{4z}\right)^{2}} ,
  \label{eq:ige}
\end{equation}
valid for $z/R \gtrsim 0.5$; equivalently, at constant thrust the induced power
is reduced by the same factor.
\end{model}

The dependence is inverse-square in $z/R$, which decides the magnitude
immediately.

\begin{example}
For $R=\SI{0.246}{\metre}$ and a rotor plane $z=\SI{\rZrotorAbove}{\metre}$
above the floor, the adopted layout of \cref{sec:results},
$(R/4z)^2 \approx 2\times10^{-3}$ and the effect is $0.2\%$. Ground effect
is negligible for a small rotor at chest height, and would only matter within
about one radius of a surface.
\end{example}

\subsection{Ceiling effect}

A ceiling above the rotor restricts the \emph{inflow} rather than the wake. The
low-pressure region above the disk is bounded, the rotor is drawn toward the
surface, and thrust rises at constant power. The functional form is taken to be
the same as \cref{eq:ige} with a separate coefficient,
\begin{equation}
  \frac{T_{\mathrm{ICE}}}{T_{\mathrm{OGE}}}\bigg|_{P}
  = \frac{1}{1-k_c\left(\dfrac{R}{4z_c}\right)^{2}} ,
  \label{eq:ice}
\end{equation}
with $z_c$ the distance to the ceiling and $k_c>1$ because at equal spacing
ceiling effect is the stronger of the two. \Cref{fig:image} shows the image
construction for both surfaces.

\begin{figure}[tbp]
\centering
% Generated by paper/figures/gen_c.py. Do not edit by hand.
\def\cM{1.724}
\def\cMbat{0.288}
\def\cL{0.3827}
\def\cR{0.2460}
\def\cReach{0.629}
\def\cSpan{1.257}
\def\cKgap{1.100}
\def\cGap{0.0492}
\def\cSoverD{0.100}
\def\cKint{1.0252}
\def\cRcp{0.130}
\def\cCoMz{0.0275}
\def\cScreenOffz{0.1025}
\def\cScreenW{0.345}
\def\cScreenH{0.194}
\def\cArmr{8.0}
\def\cWall{0.66}
\def\cWallStress{0.083}
\def\cWallStiff{0.660}
\def\cWallMin{0.60}
\def\cDriver{stiffness}
\def\cSigmaWork{44}
\def\cSigmaAllow{350}
\def\cFbend{33}
\def\cFlimit{15.2}
\def\cZrotor{1.37}
\def\cZceil{1.03}
\def\cZoverR{5.57}
\def\cZcOverR{4.19}
\def\cKground{1.0020}
\def\cKceiling{1.0072}
\def\cKc{2.0}
\def\cCeilPct{0.7}
\def\cGroundPct{0.20}
\def\cKcSpanG{13}
\def\cCeiling{2.4}
\def\cBpf{34.5}
\def\cRpm{1036}
\def\cNb{2}
\def\cAbpf{-37.5}
\def\cLwa{61.8}
\def\cRroom{20.8}
\def\cRc{0.91}
\def\cStandoff{1.20}
\def\cSplEar{56.6}
\def\cSplOne{53.8}
\def\cSplFar{54.8}
\def\cTrotor{4.23}
\def\cVtip{26.7}
\def\cAnchor{82.9}
\def\cVi{3.01}
\def\cTfCross{21}
\def\cTf{30}
\def\cQmax{0.289}
\def\cMmot{61}
\def\cKtau{5.5}
\def\cEmot{0.30}
\def\cRoomL{5}
\def\cRoomW{4}
\def\cAlphaRoom{0.20}
%
% Image-rotor construction for ground and ceiling proximity, drawn to scale
% for the design's rotor in the reference room. The mirrored wake makes the
% floor a plane of symmetry, across which no air passes.
\begin{tikzpicture}[scale=1.5, >={Stealth[length=1.6mm]}, font=\footnotesize]
  \pgfmathsetmacro{\zr}{\cZrotor}
  \pgfmathsetmacro{\H}{\cCeiling}
  \pgfmathsetmacro{\rr}{\cR}
  \pgfmathsetmacro{\zi}{2*\H-\zr}
  % floor and ceiling
  \fill[pattern=north east lines, pattern color=cG] (-1.6,-0.1) rectangle (1.6,0);
  \draw[line width=0.8pt] (-1.6,0) -- (1.6,0);
  \node[anchor=south west] at (-1.6,0.02) {floor};
  \fill[pattern=north east lines, pattern color=cG] (-1.6,\H) rectangle (1.6,\H+0.1);
  \draw[line width=0.8pt] (-1.6,\H) -- (1.6,\H);
  \node[anchor=north west] at (-1.6,\H-0.02) {ceiling};
  % the rotor, its inflow and its wake turning along the floor
  \draw[cA, line width=1.6pt] (-\rr,\zr) -- (\rr,\zr);
  \foreach \s in {-1,1} {
    \draw[cA, line width=0.6pt, ->]
      (\s*\rr,\zr) .. controls (\s*0.8*\rr,\zr-0.35) .. (\s*0.72*\rr,0.9)
      .. controls (\s*0.7*\rr,0.25) and (\s*0.6,0.07) .. (\s*1.25,0.05);
    \draw[cA, line width=0.6pt, <-]
      (\s*1.25*\rr,\zr+0.02) .. controls (\s*1.35*\rr,\zr+0.4) .. (\s*2.0*\rr,\zr+0.72);
    % the ground image and its mirrored wake
    \draw[cA!45, dashed, line width=0.6pt, ->]
      (\s*\rr,-\zr) .. controls (\s*0.8*\rr,-\zr+0.35) .. (\s*0.72*\rr,-0.9)
      .. controls (\s*0.7*\rr,-0.25) and (\s*0.6,-0.07) .. (\s*1.25,-0.05);
    % the ceiling image and its mirrored inflow
    \draw[cE!55, dashed, line width=0.6pt, <-]
      (\s*1.25*\rr,\zi-0.02) .. controls (\s*1.35*\rr,\zi-0.4) .. (\s*2.0*\rr,\zi-0.72);
  }
  \draw[cA!45, dashed, line width=1.6pt] (-\rr,-\zr) -- (\rr,-\zr);
  \draw[cE!55, dashed, line width=1.6pt] (-\rr,\zi) -- (\rr,\zi);
  \node[cA, anchor=east] at (-\rr-0.06,\zr) {rotor, radius $R$};
  \node[cA!70, anchor=east] at (-\rr-0.06,-\zr) {ground image};
  \node[cE!80, anchor=east] at (-\rr-0.06,\zi) {ceiling image};
  % dimensions on the right
  \draw[construct] (\rr+0.05,\zr) -- (1.8,\zr);
  \draw[construct] (\rr+0.05,-\zr) -- (2.1,-\zr);
  \draw[dim, black] (1.55,0) -- (1.55,\zr) node[midway, right] {$z$};
  \draw[dim, black] (1.55,\zr) -- (1.55,\H) node[midway, right] {$z_c$};
  \draw[dim, cG] (2.0,-\zr) -- (2.0,\zr) node[midway, right] {$2z$};
  \node[align=right, anchor=north east] at (-1.68,\H-0.3)
    {$z/R=\cZoverR$\\ $z_c/R=\cZcOverR$};
\end{tikzpicture}

\caption{The image construction behind \cref{eq:ige,eq:ice}, drawn to scale for
the design's rotor in the reference room. A mirror rotor below the floor, with
its wake mirrored, makes the floor a plane of symmetry across which no air
passes, which is the boundary condition the floor imposes; a mirror above the
ceiling does the same for the inflow. The images sit $2z$ and $2z_c$ away, many
rotor radii, and the corrections fall as the square of $R/4z$, which is why both
are small here.}
\label{fig:image}
\end{figure}

\begin{remark}[Honesty about $k_c$]
Published values for ceiling proximity scatter considerably across geometries
and measurement techniques~\cite{sanchezcuevas2017,conyers2019}. $k_c$ is the
least certain constant in this paper. It is exposed as a user parameter and
the engine reports the resulting factor separately so that its influence can be
bounded rather than believed. For the design here it reduces induced power by
\SI{\cCeilPct}{\percent} at $k_c=\cKc$, and the entire range $k_c\in[0.5,6]$
moves the design mass by \SI{\cKcSpanG}{\gram} (\cref{fig:surface}).
\end{remark}

\begin{figure}[tbp]
\centering
% Generated by paper/figures/gen_c.py. Do not edit by hand.
\def\cM{1.724}
\def\cMbat{0.288}
\def\cL{0.3827}
\def\cR{0.2460}
\def\cReach{0.629}
\def\cSpan{1.257}
\def\cKgap{1.100}
\def\cGap{0.0492}
\def\cSoverD{0.100}
\def\cKint{1.0252}
\def\cRcp{0.130}
\def\cCoMz{0.0275}
\def\cScreenOffz{0.1025}
\def\cScreenW{0.345}
\def\cScreenH{0.194}
\def\cArmr{8.0}
\def\cWall{0.66}
\def\cWallStress{0.083}
\def\cWallStiff{0.660}
\def\cWallMin{0.60}
\def\cDriver{stiffness}
\def\cSigmaWork{44}
\def\cSigmaAllow{350}
\def\cFbend{33}
\def\cFlimit{15.2}
\def\cZrotor{1.37}
\def\cZceil{1.03}
\def\cZoverR{5.57}
\def\cZcOverR{4.19}
\def\cKground{1.0020}
\def\cKceiling{1.0072}
\def\cKc{2.0}
\def\cCeilPct{0.7}
\def\cGroundPct{0.20}
\def\cKcSpanG{13}
\def\cCeiling{2.4}
\def\cBpf{34.5}
\def\cRpm{1036}
\def\cNb{2}
\def\cAbpf{-37.5}
\def\cLwa{61.8}
\def\cRroom{20.8}
\def\cRc{0.91}
\def\cStandoff{1.20}
\def\cSplEar{56.6}
\def\cSplOne{53.8}
\def\cSplFar{54.8}
\def\cTrotor{4.23}
\def\cVtip{26.7}
\def\cAnchor{82.9}
\def\cVi{3.01}
\def\cTfCross{21}
\def\cTf{30}
\def\cQmax{0.289}
\def\cMmot{61}
\def\cKtau{5.5}
\def\cEmot{0.30}
\def\cRoomL{5}
\def\cRoomW{4}
\def\cAlphaRoom{0.20}
%
% Fractional reduction in induced power from ground and ceiling effect,
% 1 - 1/k = k_s (R/4z)^2, against distance over radius: a slope of -2 on
% log axes, with the design's two distances marked and the spread of k_c.
\begin{tikzpicture}
  \begin{loglogaxis}[paperplot, height=5.6cm,
      xlabel={distance to the surface over rotor radius, $z/R$},
      ylabel={induced power reduction [\si{\percent}]},
      xmin=0.5, xmax=12, ymin=0.03, ymax=60,
      xtick={0.5,1,2,5,10}, xticklabels={0.5,1,2,5,10},
      ytick={0.1,1,10}, yticklabels={0.1,1,10},
      legend pos=north east]
    \addplot[name path=hi, draw=none, domain=0.5:12, samples=40, forget plot]
      {100*6/(16*x^2)};
    \addplot[name path=lo, draw=none, domain=0.5:12, samples=40, forget plot]
      {100*0.5/(16*x^2)};
    \addplot[cE!18, forget plot] fill between[of=hi and lo];
    \addplot[cB, domain=0.5:12, samples=40] {100/(16*x^2)};
    \addlegendentry{ground, $(R/4z)^{2}$}
    \addplot[cE, domain=0.5:12, samples=40] {100*\cKc/(16*x^2)};
    \addlegendentry{ceiling, $k_c(R/4z_c)^{2}$, $k_c=\cKc$}
    \addlegendimage{area legend, fill=cE!18, draw=none}
    \addlegendentry{ceiling, $k_c\in[0.5,6]$}
    \addplot[only marks, mark=*, mark size=2.2pt, cB] coordinates {(\cZoverR,\cGroundPct)};
    \node[lbl, cB, anchor=north east] at (axis cs:\cZoverR,\cGroundPct) {floor, \SI{\cGroundPct}{\percent}};
    \addplot[only marks, mark=*, mark size=2.2pt, cE] coordinates {(\cZcOverR,\cCeilPct)};
    \node[lbl, cE, anchor=south west] at (axis cs:\cZcOverR,\cCeilPct) {ceiling, \SI{\cCeilPct}{\percent}};
  \end{loglogaxis}
\end{tikzpicture}

\caption{Fractional reduction of induced power from \cref{eq:ige,eq:ice},
$1-1/k=k_s(R/4z)^2$ with $k_s=1$ for the floor and $k_s=k_c$ for the ceiling,
against distance over rotor radius. On logarithmic axes both are lines of slope
$-2$; the band spans the published range $k_c\in[0.5,6]$. At the design's two
distances the floor is worth \SI{\cGroundPct}{\percent} and the ceiling
\SI{\cCeilPct}{\percent}, and even the upper edge of the band stays below
\SI{3}{\percent}.}
\label{fig:surface}
\end{figure}

\subsection{Rotor-to-rotor interference}

Adjacent rotors each operate partly in the other's induced flow. The penalty is
negligible when the tip-to-tip gap exceeds roughly a quarter diameter and grows
as the disks approach.

\begin{model}[Interference factor]
\label{mod:interference}
With $s$ the tip-to-tip gap and $D$ the diameter,
\begin{equation}
  \kappa_{\mathrm{int}}(s/D) =
  \begin{cases}
    1, & s/D \ge 0.25,\\[2pt]
    1 + \varepsilon\left(1 - \dfrac{s/D}{0.25}\right)^{2},
      & 0 \le s/D < 0.25,
  \end{cases}
  \qquad \varepsilon = 0.07 ,
  \label{eq:interference}
\end{equation}
applied as a multiplier on induced power. The quadratic interpolation is chosen
so that the penalty and its derivative vanish at the matching point, consistent
with an induced-flow interaction that decays smoothly.
\end{model}

For a quadrotor with $k_{\mathrm{gap}}=1.10$ in \cref{eq:armlength}, the
adjacent hub spacing is $2L\sin(\pi/4)=2.2R$ and the tip-to-tip gap is
$0.2R = 0.1D$, giving $\kappa_{\mathrm{int}} = 1.025$ (\cref{fig:interference}).

\begin{figure}[tbp]
\centering
% Generated by paper/figures/gen_c.py. Do not edit by hand.
\def\cM{1.724}
\def\cMbat{0.288}
\def\cL{0.3827}
\def\cR{0.2460}
\def\cReach{0.629}
\def\cSpan{1.257}
\def\cKgap{1.100}
\def\cGap{0.0492}
\def\cSoverD{0.100}
\def\cKint{1.0252}
\def\cRcp{0.130}
\def\cCoMz{0.0275}
\def\cScreenOffz{0.1025}
\def\cScreenW{0.345}
\def\cScreenH{0.194}
\def\cArmr{8.0}
\def\cWall{0.66}
\def\cWallStress{0.083}
\def\cWallStiff{0.660}
\def\cWallMin{0.60}
\def\cDriver{stiffness}
\def\cSigmaWork{44}
\def\cSigmaAllow{350}
\def\cFbend{33}
\def\cFlimit{15.2}
\def\cZrotor{1.37}
\def\cZceil{1.03}
\def\cZoverR{5.57}
\def\cZcOverR{4.19}
\def\cKground{1.0020}
\def\cKceiling{1.0072}
\def\cKc{2.0}
\def\cCeilPct{0.7}
\def\cGroundPct{0.20}
\def\cKcSpanG{13}
\def\cCeiling{2.4}
\def\cBpf{34.5}
\def\cRpm{1036}
\def\cNb{2}
\def\cAbpf{-37.5}
\def\cLwa{61.8}
\def\cRroom{20.8}
\def\cRc{0.91}
\def\cStandoff{1.20}
\def\cSplEar{56.6}
\def\cSplOne{53.8}
\def\cSplFar{54.8}
\def\cTrotor{4.23}
\def\cVtip{26.7}
\def\cAnchor{82.9}
\def\cVi{3.01}
\def\cTfCross{21}
\def\cTf{30}
\def\cQmax{0.289}
\def\cMmot{61}
\def\cKtau{5.5}
\def\cEmot{0.30}
\def\cRoomL{5}
\def\cRoomW{4}
\def\cAlphaRoom{0.20}
%
% The interference factor of eq:interference against the tip gap over
% diameter, with the geometry of two adjacent disks and the design point.
\begin{tikzpicture}
  \begin{axis}[halfplot, width=0.56\linewidth, height=5.2cm,
      xlabel={tip-to-tip gap over diameter, $s/D$},
      ylabel={$\kappa_{\mathrm{int}}$},
      xmin=0, xmax=0.4, ymin=0.995, ymax=1.08,
      yticklabel style={/pgf/number format/fixed, /pgf/number format/precision=2}]
    \addplot[cA, domain=0:0.25, samples=60] {1 + 0.07*(1 - x/0.25)^2};
    \addplot[cA, domain=0.25:0.4, samples=2] {1};
    \draw[construct] (axis cs:0.25,0.995) -- (axis cs:0.25,1.08);
    \node[lbl, cG, anchor=south west] at (axis cs:0.25,1.035) {$s/D=0.25$};
    \node[lbl, cA, anchor=south west] at (axis cs:0.005,1.07) {$1+\varepsilon$};
    \addplot[only marks, mark=*, mark size=2.3pt, cA] coordinates {(\cSoverD,\cKint)};
    \node[lbl, cA, anchor=south west] at (axis cs:\cSoverD,\cKint) {design};
  \end{axis}
  % geometry of two adjacent rotors
  \begin{scope}[shift={(11.4,2.0)}, scale=0.78, every node/.style={scale=1}]
    \pgfmathsetmacro{\rr}{1.0}
    \pgfmathsetmacro{\ss}{0.42}
    \fill[cA!10] (-\rr-0.5*\ss,0) circle (\rr);
    \draw[cA, line width=0.8pt] (-\rr-0.5*\ss,0) circle (\rr);
    \fill[cA!10] (\rr+0.5*\ss,0) circle (\rr);
    \draw[cA, line width=0.8pt] (\rr+0.5*\ss,0) circle (\rr);
    \draw[dim, black] (-0.5*\ss,0) -- (0.5*\ss,0);
    \node[lbl, anchor=south] at (0,0.08) {$s$};
    \draw[dim, black] (-2*\rr-0.5*\ss,-1.25) -- (-0.5*\ss,-1.25) node[midway, below, lbl] {$D$};
    \draw[construct] (-1.0-0.5*\ss,0) -- (1.0+0.5*\ss,0);
    \draw[dim, cG] (-\rr-0.5*\ss,1.25) -- (\rr+0.5*\ss,1.25)
      node[midway, above, lbl, cG] {$2L\sin(\pi/N)$};
    \node[lbl, align=center, anchor=north] at (0,-1.7) {$s=2L\sin(\pi/N)-2R$};
  \end{scope}
\end{tikzpicture}

\caption{The interference factor \cref{eq:interference} against tip gap over
diameter, with the geometry that defines the gap. The penalty and its slope
vanish together at $s/D=0.25$, so the model joins the isolated-rotor case
smoothly. The design, with $k_{\mathrm{gap}}=\cKgap$ in \cref{eq:armlength},
sits at $s/D=\cSoverD$ and pays $\kappa_{\mathrm{int}}=\cKint$.}
\label{fig:interference}
\end{figure}

\subsection{Net effect and why it is worth stating}

The three corrections enter the induced power as
\begin{equation}
  \Pind = \frac{\kappa\,\kappa_{\mathrm{int}}}
               {k_{\mathrm{ground}}\,k_{\mathrm{ceiling}}}\;T\vind ,
  \label{eq:kappatotal}
\end{equation}
and they partly cancel. \Cref{tab:nearfield} gives the measured influence on a
converged design.

\begin{table}[t]
\centering
\caption{Effect of the near-field corrections on the converged design of
\cref{sec:results}; the four rows differ only in the corrections applied.
$\kappa_{\mathrm{tot}}$ is the induced power factor actually used,
\cref{eq:kappatotal}. The whole span of the four combinations is
\SI{\rNfSpanDb}{\decibel} and \SI{\rNfSpanG}{\gram}.}
\label{tab:nearfield}
\begin{tabular}{@{}lrrrr@{}}
\toprule
Model & $\kappa_{\mathrm{tot}}$ & Gross mass & Hover power & $L_p(\SI{1}{\metre})$ \\
\midrule
\inputrows{results/tab_nearfield_rows}
\end{tabular}
\end{table}

\begin{finding}
The corrections were worth adding not because they changed the design but
because they bounded their own importance. A caveat that says ``ground and
ceiling effect are not modelled'' is unfalsifiable; a table that says they are
worth a quarter of a decibel is a result. The value of a model is often in the
size of the term it lets you discard.
\end{finding}

\subsection{Recirculation: the correction that is not made}

\Cref{mod:disk} assumes the fluid far upstream is at rest. In a closed room
this is false, and the timescale over which it becomes false is computable.

\begin{definition}[Room exchange time]
The volumetric flow through the disks is $\dot{V} = A\vind$. For a room of
volume $V_{\mathrm{room}}$,
\begin{equation}
  t_{\mathrm{ex}} \coloneqq \frac{V_{\mathrm{room}}}{A\vind}
  \label{eq:exchange}
\end{equation}
is the time for a volume of air equal to the room's to pass through the rotor
disks once.
\end{definition}

\begin{example}
For the design of \cref{sec:results}, $A=\SI{0.760}{\metre\squared}$ and
$\vind=\SI{3.0}{\metre\per\second}$, so in a
$5\times4\times\SI{2.4}{\metre}$ room
$\dot V = \SI{2.28}{\metre\cubed\per\second}$,
$V_{\mathrm{room}}=\SI{48}{\metre\cubed}$ and
$t_{\mathrm{ex}} = \SI{21}{\second}$. The mean return velocity distributed over
the floor plan is $\dot V/(5\times4) = \SI{0.11}{\metre\per\second}$.
\end{example}

\Cref{fig:recirc} shows the circulation this sets up in the room.

\begin{figure}[tbp]
\centering
% Generated by paper/figures/gen_c.py. Do not edit by hand.
\def\cM{1.724}
\def\cMbat{0.288}
\def\cL{0.3827}
\def\cR{0.2460}
\def\cReach{0.629}
\def\cSpan{1.257}
\def\cKgap{1.100}
\def\cGap{0.0492}
\def\cSoverD{0.100}
\def\cKint{1.0252}
\def\cRcp{0.130}
\def\cCoMz{0.0275}
\def\cScreenOffz{0.1025}
\def\cScreenW{0.345}
\def\cScreenH{0.194}
\def\cArmr{8.0}
\def\cWall{0.66}
\def\cWallStress{0.083}
\def\cWallStiff{0.660}
\def\cWallMin{0.60}
\def\cDriver{stiffness}
\def\cSigmaWork{44}
\def\cSigmaAllow{350}
\def\cFbend{33}
\def\cFlimit{15.2}
\def\cZrotor{1.37}
\def\cZceil{1.03}
\def\cZoverR{5.57}
\def\cZcOverR{4.19}
\def\cKground{1.0020}
\def\cKceiling{1.0072}
\def\cKc{2.0}
\def\cCeilPct{0.7}
\def\cGroundPct{0.20}
\def\cKcSpanG{13}
\def\cCeiling{2.4}
\def\cBpf{34.5}
\def\cRpm{1036}
\def\cNb{2}
\def\cAbpf{-37.5}
\def\cLwa{61.8}
\def\cRroom{20.8}
\def\cRc{0.91}
\def\cStandoff{1.20}
\def\cSplEar{56.6}
\def\cSplOne{53.8}
\def\cSplFar{54.8}
\def\cTrotor{4.23}
\def\cVtip{26.7}
\def\cAnchor{82.9}
\def\cVi{3.01}
\def\cTfCross{21}
\def\cTf{30}
\def\cQmax{0.289}
\def\cMmot{61}
\def\cKtau{5.5}
\def\cEmot{0.30}
\def\cRoomL{5}
\def\cRoomW{4}
\def\cAlphaRoom{0.20}
%
% Section through the reference room: the rotor wake strikes the floor,
% spreads, rises along the walls and returns across the ceiling into the
% rotor inflow, so the inflow stops being quiescent after t_ex.
\begin{tikzpicture}[scale=2.3, >={Stealth[length=1.6mm]}, font=\footnotesize]
  \pgfmathsetmacro{\W}{\cRoomL}
  \pgfmathsetmacro{\H}{\cCeiling}
  \pgfmathsetmacro{\xm}{2.1}
  \pgfmathsetmacro{\zr}{\cZrotor}
  \draw[line width=0.9pt] (0,0) rectangle (\W,\H);
  \fill[pattern=north east lines, pattern color=cG] (0,-0.08) rectangle (\W,0);
  % the machine and the user
  \draw[cA, line width=1.6pt] (\xm-0.5,\zr) -- (\xm+0.5,\zr);
  \fill[cE!40] (\xm-0.01,\zr+0.03) rectangle (\xm+0.01,\zr+0.23);
  \draw[black!70, line width=0.8pt] (\xm+1.2,0) -- (\xm+1.2,1.25);
  \fill[black!70] (\xm+1.2,1.4) circle (0.1);
  \node[anchor=west] at (\xm+1.3,1.4) {user};
  % downwash column
  \foreach \dx in {-0.35,0,0.35}
    \draw[cA, line width=0.9pt, ->] (\xm+\dx,\zr-0.03) -- (\xm+0.85*\dx,0.25);
  % recirculating streamlines
  \draw[cA!80, line width=0.7pt, ->] (\xm-0.3,0.22) .. controls (\xm-1.2,0.05) and (0.15,0.05) .. (0.12,0.8)
     .. controls (0.1,\H-0.2) and (0.8,\H-0.1) .. (\xm-1.0,\H-0.2)
     .. controls (\xm-0.6,\H-0.3) and (\xm-0.4,\zr+0.5) .. (\xm-0.3,\zr+0.06);
  \draw[cA!80, line width=0.7pt, ->] (\xm+0.3,0.22) .. controls (\xm+1.2,0.05) and (\W-0.15,0.05) .. (\W-0.12,0.8)
     .. controls (\W-0.1,\H-0.2) and (\W-0.9,\H-0.1) .. (\xm+1.1,\H-0.2)
     .. controls (\xm+0.6,\H-0.3) and (\xm+0.4,\zr+0.5) .. (\xm+0.3,\zr+0.06);
  \draw[cA!55, line width=0.6pt, ->] (\xm-0.2,0.35) .. controls (\xm-0.8,0.25) and (0.6,0.35) .. (0.55,1.0)
     .. controls (0.55,\H-0.5) and (\xm-0.9,\H-0.55) .. (\xm-0.15,\zr+0.1);
  \draw[cA!55, line width=0.6pt, ->] (\xm+0.2,0.35) .. controls (\xm+0.8,0.25) and (\W-0.6,0.35) .. (\W-0.55,1.0)
     .. controls (\W-0.55,\H-0.5) and (\xm+0.9,\H-0.55) .. (\xm+0.15,\zr+0.1);
  % labels
  \node[anchor=north, align=center] at (0.5*\W,-0.12)
    {$\dot V=A\vind\approx\SI{2.3}{\metre\cubed\per\second}$,\quad
     $t_{\mathrm{ex}}=V_{\mathrm{room}}/\dot V\approx\SI{21}{\second}$,\quad
     mean return velocity $\dot V/(\text{floor area})\approx\SI{0.1}{\metre\per\second}$};
  \draw[dim, black] (0,\H+0.12) -- (\W,\H+0.12) node[midway, above] {\SI{\cRoomL}{\metre}};
  \draw[dim, black] (\W+0.12,0) -- (\W+0.12,\H) node[midway, right] {\SI{\cCeiling}{\metre}};
  \node[cA, anchor=south] at (\xm,\zr+0.25) {machine};
\end{tikzpicture}

\caption{Recirculation in the reference room, schematic. The wake strikes the
floor, spreads, rises along the walls and returns across the ceiling into the
rotor inflow. At $\dot V=A\vind$ a volume of air equal to the room's passes
through the disks every $t_{\mathrm{ex}}$ of \cref{eq:exchange}, so within the
first minute of a session the machine is drawing in air it has already
accelerated, and the quiescent inflow of \cref{mod:disk} no longer holds.}
\label{fig:recirc}
\end{figure}

\begin{finding}
Within one minute of a thirty-minute session the machine is flying in air it
has already accelerated. The quiescent-inflow assumption underlying every
result of \cref{sec:momentum} has quietly ceased to hold. Correcting this is
not a matter of a coefficient: it requires solving the room flow, which is a
computational fluid dynamics problem with a moving boundary. It is the largest
outstanding gap in the model and is recorded as such in
\cref{sec:limitations}. A secondary consequence is a persistent draught of
order \SI{0.1}{\metre\per\second} throughout the room, far below the
\SI{\rWake}{\metre\per\second} wake velocity but present continuously and over the
whole floor area.
\end{finding}
